\begin{englishabstract}
Self-administered questionnaires present accessibility limitations, particularly for elderly patients and those with physical impairments. Consequently, nurse-administered interviews are required, which demand substantial personnel and time. This thesis presents a Thai voice-based conversational agent for collecting electronic patient-reported outcomes (ePRO) to reduce the burden on both patients and healthcare professionals. The system supports three assessment instruments: the Disabilities of the Arm, Shoulder and Hand (DASH) questionnaire and the Toronto Extremity Salvage Score (TESS) for the upper and lower extremities. Each item is presented through pre-generated speech accompanied by an illustration; the patient's spoken response is converted to text by automatic speech recognition (ASR), and a large language model (LLM) processes the transcript and maps it to the corresponding response option of the item. All components are deployed on a local server to preserve patient privacy. Readiness for clinical use is assessed by measuring the agreement between system-derived scores and those obtained from nurse-administered interviews at both the item and total-score levels.
% TODO: summarise key results, e.g., item-level weighted kappa and total-score ICC
\end{englishabstract}
